% GENERATED FILE. Edit canonical lessons, then run npm run generate:books.
% canonical-source-hash: fnv1a64:79a9424ab66ea5c5
% canonical-lessons: RU-C263-po-moemu, RU-C263-tak-chto, RU-C263-review

\chapter{What I Think, and So}
\label{ch:ru-what-i-think-and-so}
\hlchaptermodality{263}

\begin{chapteropening}
\textbf{By the end of this chapter:} \emph{I can give my opinion with \ru{я} \ru{думаю}, \ru{что}, \ru{по}-\ru{моему} and \ru{мне} \ru{кажется}, \ru{что}, and draw a result with \ru{поэтому} and \ru{так} \ru{что}.}

Complete RU-C263-review: an opinion said three ways and one reason turned round, \ru{потому} \ru{что} into \ru{поэтому} and \ru{так} \ru{что}.
\end{chapteropening}

\section[\texorpdfstring{\ru{я} \ru{думаю}, \ru{что} … … (ya dúmayu, shto …)}{я думаю, что … … (ya dúmayu, shto …)}]{\ru{по}-\ru{моему} and \ru{мне} \ru{кажется} — an opinion, three ways}

\label{lesson:RU-C263-po-moemu}



\begin{quote}
Say \textquotedblleft{}I think that you understand\textquotedblright{}. (\emph{Ya dúmayu, shto ty ponimáyesh.}) Then the word for \textquotedblleft{}an opinion\textquotedblright{}. (\emph{Mnéniye.})

A story told with \emph{\ru{рассказать}} says what happened. Giving reasons starts one step earlier, with what you think of it.
\end{quote}

\begin{grammarlens}[title={Grammar Lens: three frames}]
\textbf{\ru{Я} \ru{думаю}, \ru{что} …} is the plain frame, and it needs \textbf{\ru{что}} and its comma.

\textbf{\ru{По}-\ru{моему}, …} — \emph{po-móyemu} — \textquotedblleft{}in my view\textquotedblright{}. It takes no \emph{\ru{что}}; a comma follows it, and then the opinion:

\begin{quote}
\textbf{\ru{По}-\ru{моему}, \ru{это} \ru{хороший} \ru{фильм}.} — \emph{po-móyemu, éto khoróshiy film} — In my view, it's a good film.
\end{quote}
\end{grammarlens}

\begin{grammarlens}[title={Grammar Lens: \ru{мне} \ru{кажется}, \ru{что} …}]
\textbf{\ru{Мне} \ru{кажется}, \ru{что} …} — \emph{mne kázhetsya, shto} — \textquotedblleft{}it seems to me that\textquotedblright{}. It is softer than \emph{\ru{я} \ru{думаю}}: less sure, and politer when you disagree.

\begin{quote}
\textbf{\ru{Мне} \ru{кажется}, \ru{что} \ru{фильм} \ru{скучный}.} — \emph{mne kázhetsya, shto film skúchnyy} — It seems to me the film is boring.
\end{quote}

\textbf{\ru{Мне}} is \textquotedblleft{}to me\textquotedblright{}, a form of \emph{\ru{я}}. The person something seems to is not the subject here: the opinion is.
\end{grammarlens}

\subsection*{The exchange}
\begin{quote}
\textbf{— \ru{По}-\ru{моему}, \ru{это} \ru{хороший} \ru{фильм}. \ru{А} \ru{ты} \ru{как} \ru{думаешь}?} \textbf{— \ru{Я} \ru{тоже} \ru{так} \ru{думаю}.}
\end{quote}

\emph{A ty kak dúmayesh?} — \textquotedblleft{}And what do you think?\textquotedblright{} \emph{Ya tózhe tak dúmayu} — \textquotedblleft{}I think so too.\textquotedblright{} To disagree gently, the second voice says instead: \textbf{\ru{А} \ru{мне} \ru{кажется}, \ru{что} \ru{он} \ru{скучный}.} \emph{\ru{Он}}, \textquotedblleft{}he\textquotedblright{}, stands for \emph{\ru{фильм}}, a masculine noun.

\begin{cousinweb}
\textbf{\ru{по}-\ru{моему}} is \emph{\ru{по}}, \textquotedblleft{}by, after\textquotedblright{}, with \emph{\ru{моему}}, a form of \textbf{\ru{мой}}, \textquotedblleft{}my\textquotedblright{}, which EtymDB 2.1 traces to Proto-Slavic \emph{*mojĭ}: \textquotedblleft{}after my fashion\textquotedblright{}. \textbf{\ru{кажется}} is \emph{\ru{казаться}}, \textquotedblleft{}to show oneself, to seem\textquotedblright{}, the \emph{-\ru{ся}} form of an older \emph{\ru{казать}}, \textquotedblleft{}to show\textquotedblright{}. So \emph{\ru{мне} \ru{кажется}} says \textquotedblleft{}it shows itself to me\textquotedblright{}.
\end{cousinweb}

\subsection*{Your turn}
\begin{itemize}
\raggedright
  \item \emph{Say it:} \emph{Po-móyemu, éto khoróshiy film.}
  \item \emph{Say it:} \emph{Mne kázhetsya, shto film skúchnyy.}
  \item \emph{Say it:} one opinion of yours three ways: \emph{ya dúmayu, shto}; \emph{po-móyemu}; \emph{mne kázhetsya, shto}
  \item \emph{Answer:} \emph{A ty kak dúmayesh?}
\end{itemize}

\begin{tcolorbox}[breakable,colback=teal!4,colframe=teal!35!black,title={Before you move on}]
\textquotedblleft{}In my view\textquotedblright{}: which frame needs no \emph{\ru{что}}? (\emph{\textbf{\ru{По}-\ru{моему}.}}) The softer one, for \textquotedblleft{}it seems to me\textquotedblright{}? (\emph{\textbf{\ru{Мне} \ru{кажется}, \ru{что} …}}) \textquotedblleft{}I think so too\textquotedblright{}? (\emph{\textbf{\ru{Я} \ru{тоже} \ru{так} \ru{думаю}.}})
\end{tcolorbox}
\section[\texorpdfstring{\ru{поэтому}, \ru{так} \ru{что} (poétomu, tak shto)}{поэтому, так что (poétomu, tak shto)}]{\ru{поэтому} and \ru{так} \ru{что} — the reason, turned round}

\label{lesson:RU-C263-tak-chto}



\begin{quote}
Say \textquotedblleft{}because\textquotedblright{} and \textquotedblleft{}since\textquotedblright{}. (\emph{Potomú shto. Tak kak.}) Then an opinion with \emph{po-móyemu}. An opinion needs a reason, and a reason has a direction.
\end{quote}

\begin{grammarlens}[title={Grammar Lens: back to the cause, forward to the result}]
One rainy day, said two ways:

\begin{quote}
\textbf{\ru{Мы} \ru{не} \ru{гуляли}, \ru{потому} \ru{что} \ru{шёл} \ru{дождь}.} — \emph{my ne gulyáli, potomú shto shol dozhd} — We didn't go for a walk, because it was raining. \textbf{\ru{Шёл} \ru{дождь}, \ru{поэтому} \ru{мы} \ru{не} \ru{гуляли}.} — \emph{shol dozhd, poétomu my ne gulyáli} — It was raining, so we didn't go for a walk.
\end{quote}

\textbf{\ru{Потому} \ru{что}} stands in front of the cause. \textbf{\ru{Поэтому}} stands in front of the result. The facts are the same; the order flips, and the comma moves with the seam. \textbf{\ru{Поэтому}} can also open a new sentence: \emph{\ru{Шёл} \ru{дождь}. \ru{Поэтому} \ru{мы} \ru{не} \ru{гуляли}.}
\end{grammarlens}

\subsection*{You'll want to know: \ru{так} \ru{что}}
\textbf{\ru{так} \ru{что}} — \emph{tak shto} — \textquotedblleft{}so, so then\textquotedblright{}. It also points to the result, and it sounds more like talk:

\begin{quote}
\textbf{\ru{Шёл} \ru{дождь}, \ru{так} \ru{что} \ru{мы} \ru{сидели} \ru{в} \ru{кафе}.} — \emph{shol dozhd, tak shto my sidéli v kafé} — It was raining, so we sat in a café.
\end{quote}

It usually sits inside one sentence, after a comma.

\begin{cousinweb}
\textbf{\ru{так}}, \textquotedblleft{}in that way, so\textquotedblright{}, goes back to Old East Slavic \emph{tako} (EtymDB 2.1), on the pointing stem \emph{*to-} that also gives \emph{\ru{тот}}, \textquotedblleft{}that\textquotedblright{}. So \textbf{\ru{так} \ru{что}} is \textquotedblleft{}that way, so that\textquotedblright{}. \textbf{\ru{Поэтому}} is \emph{\ru{по}}, \textquotedblleft{}by\textquotedblright{}, with \emph{\ru{этому}}, a form of \emph{\ru{этот}}: \textquotedblleft{}for this\textquotedblright{}.
\end{cousinweb}

\subsection*{Your turn}
Say these aloud:

\begin{itemize}
\raggedright
  \item \emph{My ne gulyáli, potomú shto shol dozhd.}
  \item turned round: \emph{Shol dozhd, poétomu my ne gulyáli.}
  \item \emph{Shol dozhd, tak shto my sidéli v kafé.}
  \item a reason of your own, first with \emph{potomú shto}, then turned round with \emph{poétomu}
\end{itemize}

\begin{tcolorbox}[breakable,colback=teal!4,colframe=teal!35!black,title={Before you move on}]
Which word stands in front of the cause? (\emph{\textbf{\ru{Потому} \ru{что}.}}) And in front of the result? (\emph{\textbf{\ru{Поэтому}}}, or \emph{\textbf{\ru{так} \ru{что}}}.) Which of the two opens a new sentence more easily? (\emph{\textbf{\ru{Поэтому}.}})
\end{tcolorbox}
\section[(review)]{Review — what I think, and so}

\label{lesson:RU-C263-review}



\begin{quote}
Say \textquotedblleft{}in my view, it's a good film\textquotedblright{}. (\emph{Po-móyemu, éto khoróshiy film.}) Then softer: \textquotedblleft{}it seems to me the film is boring\textquotedblright{}. (\emph{Mne kázhetsya, shto film skúchnyy.})
\end{quote}

\begin{grammarlens}[title={What you've built}]
Two lessons, and an opinion now comes with its reason. Three frames: \textbf{\ru{я} \ru{думаю}, \ru{что}}, \textbf{\ru{по}-\ru{моему}}, \textbf{\ru{мне} \ru{кажется}, \ru{что}}. And two directions: \textbf{\ru{потому} \ru{что}} before the cause, \textbf{\ru{поэтому}} or \textbf{\ru{так} \ru{что}} before the result.

All of it in two sentences:

\begin{quote}
\textbf{\ru{По}-\ru{моему}, \ru{музей} \ru{интересный}, \ru{потому} \ru{что} \ru{он} \ru{старый}.} — \emph{po-móyemu, muzéy interésnyy, potomú shto on stáryy} — In my view the museum is interesting, because it's old. \textbf{\ru{Музей} \ru{старый}, \ru{поэтому}, \ru{мне} \ru{кажется}, \ru{он} \ru{интересный}.} — the same reason, pointed forward, and said more softly.
\end{quote}
\end{grammarlens}

\subsection*{Your turn}
\begin{itemize}
\raggedright
  \item \emph{Say it:} \emph{Po-móyemu, muzéy interésnyy, potomú shto on stáryy.}
  \item \emph{Say it:} \emph{Shol dozhd, tak shto my sidéli v kafé.}
  \item \emph{Answer:} \emph{A ty kak dúmayesh?} about a film you know, with \emph{mne kázhetsya, shto}
  \item \emph{Say it:} a short Saturday story: a scene, then \emph{snachála}, \emph{zatém}, and a result with \emph{poétomu}
\end{itemize}

\begin{tcolorbox}[breakable,colback=teal!4,colframe=teal!35!black,title={Before you move on}]
\textquotedblleft{}So\textquotedblright{}, in front of a result, two ways? (\emph{\textbf{\ru{Поэтому}}}, \emph{\textbf{\ru{так} \ru{что}}}.) \textquotedblleft{}In my view\textquotedblright{}, with no \emph{\ru{что}}? (\emph{\textbf{\ru{По}-\ru{моему}.}})
\end{tcolorbox}
